\section{Codes for the Synthesizer}
\label{app:synth_Code}

\subsection{Oscillator}

The following header file below shows how the basic waveforms shapes are defined, which only consists
of the four basic waveforms: Sine, Square, Sawtooth and Triangle. The samples depend on the location of the
phase, between 0 and 1.

\lstinputlisting[
    style=cppstyle,
    caption={Header file for defining the Oscillator shapes},
    label={lst:oscshape}
]{../include/audio_source/oscillator/OscShape.hpp}

The following listing below shows how the program renders the oscillator signal and sent to the audio engine.
The callback function is stateful, so everytime it runs, the phase variable moves on and recalculates the next
sample \cite{openai_chatgpt}.

\lstinputlisting[
    style=cppstyle,
    linerange={osc_generate-osc_generate},
    caption={Stateful callback function for generating waveforms},
    label={lst:osc_generate}
]{../include/audio_source/oscillator/Oscillator.cpp}

\subsection{State Variable Filter}
\label{app:synth_code}

The following listing show how the program updates parameters for resonance ($R = \frac{1}{2Q}$) and the
coefficient, $g$ \eqref{eq:zdfsvf_g}. They perform only a simple arithmetic which potentially reduces
computational cost to a greater extent compared to the Biquad.

\lstinputlisting[
    style=cppstyle,
    linerange={svf_setters-svf_setters},
    caption={Stateful callback function for generating waveforms},
    label={lst:svf_setters}
]{../include/audio_fx/eq_filter/StateVariableFilter.cpp}

This listing below represents Zavalshin's algorithm for the ZDF-SVF filter
\cite{vafilterdesign2021}.

\lstinputlisting[
    style=cppstyle,
    linerange={svf_process-svf_process},
    caption={Stateful callback function for the ZDF-SVF Filter},
    label={lst:svf_process}
]{../include/audio_fx/eq_filter/StateVariableFilter.cpp}

\subsection{Creating Unison}
\label{app:unison_code}

The following listing below shows a complete header file for forming unison from an 
oscillator, as described in Section \ref{sec:unison}.

\lstinputlisting[
    style=cppstyle,
    % linerange={svf_process-svf_process},
    caption={Class Structure for Creating Unison},
    label={lst:unison}
]{../include/track/instrument/osc_synth/Unison.hpp}
